\section{Introduction}
The rapid advancement of large language models (LLMs) has driven a paradigm shift from passive text generators to active, \emph{API-calling agents}~\citep{schick2023toolformer, yao2023react}. These agents fulfill complex user requests by executing API calls and reasoning over environmental feedback. However, training such agents requires vast amounts of high-quality trajectory data that capture the full multi-step interaction process, including user instructions, sequences of interleaved read and write API calls, execution results, and final responses. Compared to standard instruction-tuning data, these multi-step trajectories are substantially harder to collect at scale, yet they remain essential for teaching agents robust API usage, accurate state tracking, and effective error recovery.

Given the difficulty of manual collection, synthetic data generation has emerged as a natural solution, with recent works leveraging LLMs to produce agentic trajectories at scale~\citep{toolllm, agenttuning, apigen, xu2026controllable, lu2026firefly}. However, these works assume access to fully implemented environments with executable APIs and backend databases filled with realistic underlying states. This tight coupling between environment creation and data synthesis creates a major scalability bottleneck. To generate training data for a new set of APIs, one must first engineer the corresponding infrastructure and, more challengingly, source the diverse underlying state configurations required to populate the environment.

\begin{figure*}[t]
    \centering
    \resizebox{\textwidth}{!}{%
    \begin{tikzpicture}[
        node distance=0.4cm and 0.6cm,
        >={Stealth[length=3mm]},
        every node/.style={font=\normalsize},
        stagebox/.style={
            rounded corners=4pt,
            minimum height=1.4cm,
            minimum width=3.6cm,
            align=center,
            text=white,
            font=\normalsize\bfseries,
            drop shadow={shadow xshift=0.5mm, shadow yshift=-0.5mm, opacity=0.3},
        },
        subbox/.style={
            rounded corners=3pt,
            minimum height=0.85cm,
            minimum width=3.0cm,
            align=center,
            font=\small,
            draw=gray!50,
            fill=white,
        },
        outputbox/.style={
            rounded corners=4pt,
            minimum height=1.4cm,
            minimum width=3.2cm,
            align=center,
            text=white,
            font=\normalsize\bfseries,
            fill=gray!55!black,
            drop shadow={shadow xshift=0.5mm, shadow yshift=-0.5mm, opacity=0.3},
        },
        docfield/.style={
            font=\small\ttfamily,
            anchor=west,
            inner sep=1pt,
        },
        doclabel/.style={
            font=\small\bfseries,
            text=gray!50!black,
            anchor=west,
            inner sep=1pt,
        },
        arrowlabel/.style={
            font=\normalsize,
            text=gray!70!black,
        },
        innerarrow/.style={
            ->, thick, gray!60!black, shorten >=1pt, shorten <=1pt,
        },
    ]

    % ---- Stage 1 ----
    \node[stagebox, fill=stage1] (s1) {Task Synthesis};
    \node[subbox, below=0.45cm of s1] (s1a) {Task generator LLM\\ composes tasks};
    \node[subbox, below=0.45cm of s1a] (s1b) {Judge LLM filters\\ out invalid tasks};
    \draw[innerarrow] (s1a.south) -- (s1b.north);

    % ---- Stage 2 ----
    \node[stagebox, fill=stage2, right=3.0cm of s1] (s2) {Trajectory Synthesis};
    \node[subbox, below=0.45cm of s2] (s2a) {Agent LLM issues\\API calls};
    \node[subbox, below=0.45cm of s2a] (s2b) {Simulator LLM\\synthesizes responses};

    % ---- Stage 2 arrows (labels centered on the vertical segments) ----
    \draw[->, thick, stage2!80!black] (s2a.west) -- ++(-0.3,0) |- (s2b.west);
    \node[arrowlabel, anchor=east] at ([xshift=-0.3cm]$(s2a.west)!0.5!(s2b.west)$) {API calls};

    \draw[->, thick, stage2!80!black] (s2b.east) -- ++(0.3,0) |- (s2a.east);
    \node[arrowlabel, anchor=west] at ([xshift=0.3cm]$(s2a.east)!0.5!(s2b.east)$) {responses};

    % ---- Stage 3 ----
    \node[stagebox, fill=stage3, right=3.0cm of s2] (s3) {Trajectory Filtering};
    \node[subbox, below=0.45cm of s3] (s3a) {Judge LLM evaluates\\ trajectory quality};
    \node[subbox, below=0.45cm of s3a] (s3b) {High-quality\\ trajectories selected};
    \draw[innerarrow] (s3a.south) -- (s3b.north);

    \begin{scope}[on background layer]
        \node[fit=(s1)(s1a)(s1b), rounded corners=6pt, fill=stage1bg, inner sep=6pt] (s1bg) {};
        \node[fit=(s2)(s2a)(s2b), rounded corners=6pt, fill=stage2bg, inner sep=6pt] (s2bg) {};
        \node[fit=(s3)(s3a)(s3b), rounded corners=6pt, fill=stage3bg, inner sep=6pt] (s3bg) {};
    \end{scope}

    % ---- Output ----
    \node[outputbox, anchor=south west] (output) at ([xshift=2.0cm]s3bg.south east) {Trajectory dataset\\ for training};

    % ---- Doc string card ----
    \node[
        draw=gray!60,
        rounded corners=3pt,
        fill=white,
        inner sep=6pt,
        anchor=east,
        align=left,
        label={[font=\normalsize\bfseries, text=gray!40!black, yshift=-2pt]above:API specifications},
    ] (input) at ([xshift=-2.0cm]s1bg.west) {%
        \begin{tikzpicture}[inner sep=0pt]
            \node[font=\small\ttfamily\bfseries, text=stage1!60!black, anchor=west] (name) at (0,0) {send\_money};
            \node[font=\small\itshape, text=gray!60!black, anchor=west, below=2pt of name.west] (desc) {Send money to a user.};
            \draw[gray!40, dashed] ([yshift=-5pt]desc.south west) -- ++(4.2,0);
            \node[doclabel, below=10pt of desc.south west, anchor=north west] (inlabel) {Input:};
            \node[docfield, below=1pt of inlabel.south west, anchor=north west, text=gray!30!black] (in1)
                {receiver\_id: \textcolor{typecolor}{int}};
            \node[docfield, below=0pt of in1.south west, anchor=north west, text=gray!30!black] (in2)
                {amount: \textcolor{typecolor}{float}};
            \node[docfield, below=0pt of in2.south west, anchor=north west, text=gray!30!black] (in3)
                {note: \textcolor{typecolor}{str} = \textcolor{stringcolor}{""}};
            \node[doclabel, below=3pt of in3.south west, anchor=north west] (outlabel) {Output:};
            \node[docfield, below=1pt of outlabel.south west, anchor=north west, text=gray!30!black] (out1) {\{};
            \node[docfield, below=0pt of out1.south west, anchor=north west, text=gray!30!black, xshift=10pt] (out2)
                {"transaction\_id": \textcolor{typecolor}{int},};
            \node[docfield, below=0pt of out2.south west, anchor=north west, text=gray!30!black] (out3)
                {"status": \textcolor{typecolor}{str}};
            \node[docfield, below=0pt of out3.south west, anchor=north west, text=gray!30!black, xshift=-10pt] (out4) {\}};
        \end{tikzpicture}%
    };

    % ---- Main horizontal arrows ----
    \draw[->, thick, gray!70!black] (input.east|-s1) -- (s1bg.west|-s1)
        node[arrowlabel, midway, above] {input};

    \draw[->, thick, gray!70!black] (s1bg.east|-s1) -- (s2bg.west|-s2)
        node[arrowlabel, midway, above] {valid tasks};
    \draw[->, thick, gray!70!black] (s2bg.east|-s2) -- (s3bg.west|-s3)
        node[arrowlabel, midway, above] {trajectories};

    % L-shaped arrow to output
    \draw[->, thick, gray!70!black] (s3bg.east|-s3) -| (output.north)
        node[arrowlabel, pos=0.25, above] {output};

    \end{tikzpicture}
    }%
    % \vspace{-0.05in}
    \caption{\textsc{ESAT} pipeline that generates agentic trajectories using only API specifications as input. First, a task generator LLM composes tasks based on API specifications which are then filtered by a judge LLM. Then, an agent LLM generates solution trajectories by interacting with an API simulator LLM. Finally, a trajectory judge LLM selects high quality trajectories.}
    \label{fig:pipeline}
    \vspace{-0.15in}
\end{figure*}


In this work, we ask a fundamental question: \emph{can we generate high-quality agentic trajectory data without building the underlying environments at all}? Our key conjecture is that modern LLMs, having been pretrained on vast corpora of world knowledge, can function as dynamic digital world models. Given API specifications, input arguments, task instructions, and prior interaction history as context, an LLM can synthesize realistic execution results that mimic feedback from a live environment, accurately reflecting the underlying state changes. Driven by this insight, we introduce a highly scalable paradigm for agentic data generation. Instead of constructing real environments and executing actual APIs, we use LLMs to simulate all environmental feedback on-the-fly directly from API specifications.

To instantiate this paradigm, we propose \textsc{ESAT}, an \emph{\textbf{E}nvironment-Free \textbf{S}ynthetic \textbf{A}gentic \textbf{T}rajectory} generation pipeline that uses only API specifications as input. The pipeline operates in three stages, as illustrated in \Cref{fig:pipeline}. First, during the \textit{task synthesis} stage, an LLM generates diverse, realistic tasks solvable by the target APIs, with an LLM-based critic filtering out invalid tasks. Second, the core \textit{trajectory synthesis} stage pairs a teacher agent with our LLM-based API simulator. As the teacher iteratively issues API calls to solve the task, the simulator seamlessly generates coherent environmental feedback on-the-fly using the user task and prior simulation history as context. This interaction produces complete, multi-step agent trajectories entirely in simulation, without ever requiring an executable environment (see \Cref{app:trajectories} for example trajectories). Finally, during \textit{trajectory filtering}, an LLM judge evaluates the generated trajectories for correctness, completeness, and redundancy, discarding sub-optimal samples to ensure a high-fidelity training dataset.

\textsc{ESAT} offers two critical advantages. First, it turns agentic trajectory synthesis into a lightweight specification problem. By relying solely on API specifications rather than fully implemented environments or real-world deployed services, it allows data generation to scale seamlessly across diverse API ecosystems. Second, leveraging the long-context capabilities of modern LLMs, our simulator maintains state coherence across successive API calls, allowing for the generation of rich, multi-turn training trajectories that go far beyond isolated API-calling examples.

To validate \textsc{ESAT}, we conduct extensive experiments on AppWorld~\citep{appworld} and OfficeBench \citep{wang2024officebench}, two challenging multi-app benchmarks that evaluate agents on complex, multi-step information-retrieval and state-changing tasks. Our experiments across models ranging from 1.7B to 27B parameters demonstrate that supervised fine-tuning (SFT) on \textsc{ESAT} data leads to significant performance improvements with respect to the corresponding base models, yielding up to 50.5\% on AppWorld and 60.5\% on OfficeBench. Most notably, on AppWorld, which provides an official set of training tasks, models trained on \textsc{ESAT} data outperform the models trained on the training task trajectories collected from the real AppWorld environment.

Our contributions are summarized below:
\begin{itemize}[itemsep=1mm,parsep=1pt,topsep=2pt,leftmargin=*]  
\item We introduce a context-aware, LLM-based API simulator that acts as a dynamic digital world model. By leveraging task context and interaction history, this simulator generates coherent environmental feedback, enabling the synthesis of complex, multi-step, API-calling trajectories without any real API execution.
\item We propose \textsc{ESAT}, a novel, environment-free data synthesis pipeline that reduces agentic trajectory generation to a lightweight specification problem. By requiring only API specifications as input, \textsc{ESAT} eliminates the infrastructure and data bottlenecks associated with building and populating executable environments, enabling seamless scaling across diverse API ecosystems.
\item We demonstrate that \textsc{ESAT} data is highly effective for supervised fine-tuning. Across models ranging from 1.7B to 27B parameters, our synthetic data drives substantial gains on AppWorld and OfficeBench datasets.
\end{itemize}